\documentclass{beamer}
\usecolortheme{solarized}
\usepackage{hyperref}
\setbeamertemplate{section in toc}[sections numbered]

\title{Git In Here}
\subtitle{An Unoffical Introduction To Git Version Control Management}
\author{Sam Dixon\\
{\tt \href{mailto:s.dixon@napier.ac.uk}{s.dixon@napier.ac.uk}}}
\date{2018}
\titlegraphic{\includegraphics[height=3.5cm]{../res/git_icon_orange.png}} 

\begin{document}

\frame{\titlepage}

\begin{frame}{Contents}
    \tableofcontents
\end{frame}

\section{What is Version Control?}
\begin{frame}[allowframebreaks]{\secname}
    Version Control Systems are tools that let you track changes

    A number of different VCSs are available:
    \begin{itemize}
        \item Subversion
        \item Mercurial
        \item CVS
        \item Fossil
        \item Git
    \end{itemize}
    Each VCS has its own quirks in style and implementation

\framebreak
    
    Common features of VCSs include:
    \begin{itemize}
        \item Timeline of changes
        \item File reversion  
        \item Branching
        \item Tagging
        \item Collaboration
    \end{itemize}
\end{frame}


\begin{frame}{Why you should use Version control}
    VCSs can help with common problems like:
    \begin{itemize}
        \item Creating multiple copies files ``in case you need the old one''
        \item Put work onto a USB drive to move it between computers
        \item Copied code into email/IM to send it to a friend
        \item Lost work due to loss of power/faulty hard drive
        \item Overwritten someone else's work when collaborating
    \end{itemize}
\end{frame}

\section{What is Git?}
\begin{frame}[allowframebreaks]{\secname}
    Git was developed in 2005 by Linus Torvalds.

    It has become one of the most popular VCSs in use today

    Some companies that use Git include:
    \begin{itemize}
        \item Google
        \item Apple
        \item Microsoft
        \item Facebook
        \item Amazon
        \item Adobe
        \item Mozilla
        \item NASA
    \end{itemize}

\framebreak
    
    Git has a number of key features that make it an attractive VCS
    \begin{itemize}
        \item Distributed
        \item Decentralised
        \item Free (like beer {\it and} like speech)
        \item Multiplatform
        \item Speed
        \item Popularity
    \end{itemize}
\end{frame}

\begin{frame}{Common Terms}
    Git has a number of key terms that you should keep in mind
    \begin{itemize}
        \item Repository 
        \item Commit
        \item Push/Pull
        \item Tag
        \item Branch
        \item Merge
    \end{itemize}
    Some of terms may be used in other VCSs, but their functions may not be
    the same
\end{frame}

\section{Setting up Git}
\begin{frame}[allowframebreaks]{\secname}
    Git can be installed in a number of ways:
    \begin{itemize}
        \item Download directly: {\tt \url{https://git-scm.com/downloads}}
        \item Install a Git visualiser (GitHub Desktop, Sourcetree, Kraken)
        \item Install via a package manager 
    \end{itemize}

\framebreak
    
    Git requires a name and email to create commits \\~\\
    {\tt \$ git congif --global user.name "Sam"} \\~\\
    {\tt \$ git config --global user.email sam@mail.com} \\~\\
    These settings are not used as login credentials / tokens
\end{frame}

\section{Git Workflow}
\begin{frame}{\secname}
    The standard process for using Git is as follows: \\~\\
    \begin{enumerate}
        \item {\bf Initialise} a Git repo
        \item Do some work (edit files)
        \item {\bf Stage} the changes you want to keep
        \item Finalise the changes in a {\bf Commit}
        \item Repeat from step 2
    \end{enumerate}
\end{frame}

\section{Basic Commands}
\begin{frame}{Getting a Repo}
    Git can create a repo in  a new folder:\\~\\
    {\tt \$ git init project} \\~\\
    or in an existing directory: \\~\\
    {\tt \$ git init} \\~\\
    It doesn't matter if there are files already in the directory
\end{frame}

\begin{frame}{Staging Files}
    Git will not track files by default - you have to specify: \\~\\
    {\tt \$ git add file\_name} \\~\\
    You can add multiple files in a single command: \\~\\
    {\tt \$ git add file\_1 file\_2} \\~\\
    Or you can use glob wildcards: \\~\\
    {\tt \$ git add *.txt}
\end{frame}

\begin{frame}{Checking File State}
    You can check the state of a repo with the following: \\~\\
    {\tt \$ git status} \\~\\
    Files can be in multiple states:
    \begin{itemize}
        \item Untracked
        \item Modified
        \item Staged
        \item Deleted
    \end{itemize}
\end{frame}

\begin{frame}{Committing Changes}
    When you are ready to finalise your changes; commit: \\~\\
    {\tt \$ git commit} \\~\\
    You can also declare a commit message in the same command: \\~\\
    {\tt \$ git commit -m "Your message here"}
\end{frame}

\begin{frame}{Viewing History}
    Once you have a few commits, you can review them:\\~\\
    {\tt \$ git log} \\~\\
    View what changes were made in each commit with the {\tt patch} flag: \\~\\
    {\tt \$ git log -p} \\~\\
    There are multiple flags that can alter how the log looks: \\~\\
    {\tt \$ git log --oneline --graph --decorate}
\end{frame}

\section{Undoing}
\begin{frame}{\secname}
    The {\tt amend} flag allows you to alter a commit

    It adds your staged files to the {\it previous} commit 

    Use this to add files you forgot, or for quick alterations:\\~\\
    {\tt \$ git commit\\
    \$ git add forgotten\_file\\
    \$ git commit --amend} \\~\\

    When you amend a commit you can also change it's message
\end{frame}

\begin{frame}{Unstaging a File}
    You may accidentally stage a file that you do not want to commit

    You can unstage the file, without losing the changes you made:\\~\\
    {\tt \$ git reset HEAD file\_name} \\~\\
    {\tt HEAD} is Git's name for the tip of the current branch   
\end{frame}

\begin{frame}{Unmodifying a File}
    You may edit a file and decide to not keep the modifications

    You revert a file to how it was when it was last committed: \\~\\
    
    {\tt \$ git checkout -- file\_name} \\~\\

    This command is dangerous - you cant recover uncommitted changes
\end{frame}


\section{Branching}
\begin{frame}{\secname}
    Branches are a key feature of Git

    They allow different version of the same project to exists simultaneously
    
    Allows for parallel development of features

    Branches can be used to experiment or write patches

    Branches are inexpensive to create/delete
\end{frame}

\begin{frame}{Creating Branches}
You can quickly create branches:

{\tt \$ git branch branch\_name}

This will create a new branch starting at your current commit\\~\\

You can list the different branches in a repo:

{\$ git branch}

The {\tt status} command will inform you of which branch you are on
\end{frame}

\begin{frame}{Swapping branches}
Once a branch has been created, you need to switch to it:

{\tt \$ git checkout branch\_name} 

Git may stop you from checking out if you have uncommitted changes

\end{frame}

\begin{frame}{Merging Branches}
If you have commits on one branch and want to copy them to another you will
need to merge 

Merging will make Git apply the divergent commits from one branch to another

First, swap to the branch you want to apply the commits to:
{\$ git checkout master}

\end{frame}


\section{Remotes}
\begin{frame}[allowframebreaks]{\secname}
    Git can be used entirely locally\\~\\

    But remotes allow you to collaborate and store projects\\~\\

    A remote is copy of a Git repo ``somewhere else''\\~\\

    A number of services will host repos for you\\~\\

    You can even run your own Git server\\~\\

\framebreak

    There a two ways to set up repo remotes \\~\\

    You can make a local copy of a remote repo: \\~\\

    {\tt \$ git clone remote\_url} \\~\\

    This will copy the  repo and set the url as the default remote
    \\~\\

    You can also add a remote to a local repo manually: \\~\\

    {\tt \$ git remote add remote\_name remote\_url} \\~\\

    The default name for a repo's primary remote is {\tt origin} - but you can
    make it what ever you want

\framebreak
    
    You can see what remotes are configured to a repo: \\~\\

    {\tt \$ git remote} \\~\\

    To see the url's of theses remotes include the {\tt verbose} flag: \\~\\

    {\tt \$ git remote -v} \\~\\

    You can rename remote:\\~\\
    {\tt \$ git remote rename remote\_name new\_name} \\~\\

    Or remove it if you dont need to use it any more: \\~\\

    {\tt \$ git remote remove remote\_name}

\framebreak

    You can share local commits to a remote: \\~\\

    {\tt \$ git push remote\_name branch\_name} \\~\\

    The default branch name is {\tt master} \\~\\

    You can query a remote to learn its state: \\~\\

    {\tt \$ git remote show remote\_name} \\~\\

\framebreak

    To get commits from a remote repo, you must fetch it:\\~\\
    {\tt \$ git fetch remote\_name} \\~\\
    Fetching will copy data - but will not change any files
    
    You can choose to add the remote commits if you want
\end{frame}

\section{Tagging}
\begin{frame}[allowframebreaks]{\secname}
    Tags are simply a pointer to a specific commit

    Used to show commits of interest, like version releases or bugs \\~\\
    Create a lightweight tag with the following command: \\~\\
    {\tt \$ git tag tag\_name} \\~\\
    Annotated tags can be created as well:  \\~\\
    {\tt \$ git tag -a tag\_name -m "tag message"} \\~\\
\framebreak
    You can see what tags there are in a repo: \\~\\
    {\tt \$ git tag} \\~\\
    Or you can inspect a specific tag: \\~\\
    {\tt \$ git show tag\_name}
\end{frame}



\end{document}

